\documentclass[aspectratio=169]{beamer}

\usetheme{Madrid}
\usecolortheme{default}
\usepackage{amsmath,amssymb,mathtools}
\usepackage{bm}
\usepackage{tikz}
\usepackage{graphicx}

\title{Vectors, Vector Spaces and Linear Algebra}
\author{Gayathri P S}
\date{\today}

\begin{document}

\frame{\titlepage}

\begin{frame}{Contents}
\tableofcontents
\end{frame}


\section{Foundations}

\begin{frame}{Sets and Algebraic Structures}
\textbf{Set:} A well-defined collection of objects.
\[
A=\{1,2,3,4\}
\]
Important operations:
\[
A\cup B,\quad A\cap B,\quad A-B,\quad A^c,\quad A\times B
\]

An algebraic structure is a set with a binary operation:
\[
(S,*)
\]

\begin{center}
\begin{tabular}{|c|c|}
\hline
\textbf{Structure} & \textbf{Additional property}\\
\hline
Magma & Closure\\
Semigroup & Closure + Associativity\\
Monoid & Semigroup + Identity\\
Group & Monoid + Inverse\\
Abelian Group & Group + Commutativity\\
\hline
\end{tabular}
\end{center}

\textbf{ML connection:} Sets describe datasets/classes, while algebraic structures provide the mathematical foundation for transformations and operations.
\end{frame}

\begin{frame}{Vectors and Vector Spaces}
A vector is an ordered collection of numbers:
\[
\mathbf v=
\begin{bmatrix}
v_1\\v_2\\\vdots\\v_n
\end{bmatrix},
\qquad
\|\mathbf v\|=\sqrt{\sum_i v_i^2}.
\]

Vector addition:
\[
\mathbf u+\mathbf v
=
\begin{bmatrix}
u_1+v_1\\u_2+v_2\\\vdots
\end{bmatrix}
\]

Scalar multiplication:
\[
\alpha\mathbf v=
\begin{bmatrix}
\alpha v_1\\\alpha v_2\\\vdots
\end{bmatrix}
\]

A vector space is closed under vector addition and scalar multiplication.

\textbf{Polynomial example:}
\[
p(x)=a+bx+cx^2
\longleftrightarrow
\begin{bmatrix}a\\b\\c\end{bmatrix}
\]

\textbf{ML:} A data point is represented by a feature vector.
\end{frame}

\begin{frame}{Linear Combinations, Hyperplanes and Perceptron}
A linear combination is
\[
c_1\mathbf v_1+c_2\mathbf v_2+\cdots+c_k\mathbf v_k.
\]

The equation
\[
ax+by+c=0
\]
represents a line, while
\[
\mathbf w^T\mathbf x+b=0
\]
represents a hyperplane.

\[
\boxed{\mathbf w=\text{normal vector to the hyperplane}}
\]

A perceptron computes
\[
z=\mathbf w^T\mathbf x+b.
\]

Using
\[
\mathbf w^T\mathbf x
=
\|\mathbf w\|\|\mathbf x\|\cos\theta,
\]
the angle between $\mathbf w$ and $\mathbf x$ influences classification.

\[
\theta<90^\circ\Rightarrow z>0,\qquad
\theta>90^\circ\Rightarrow z<0
\]

Thus, a perceptron is a geometric linear classifier.
\end{frame}

\begin{frame}{Mathematical Spaces}
Different mathematical spaces provide different ways to study mathematical objects.

\begin{center}
\begin{tabular}{|c|c|}
\hline
\textbf{Space} & \textbf{Main concept}\\
\hline
Vector Space & Addition and scalar multiplication\\
Metric Space & Distance\\
Inner Product Space & Angle and orthogonality\\
Hilbert Space & Complete inner product space\\
Topological Space & Continuity and convergence\\
\hline
\end{tabular}
\end{center}

For an inner product:
\[
\langle\mathbf u,\mathbf v\rangle=\mathbf u^T\mathbf v
\]

Angle:
\[
\cos\theta=
\frac{\mathbf u^T\mathbf v}
{\|\mathbf u\|\|\mathbf v\|}
\]

\textbf{AI/ML applications:}
\[
\text{Metric}\rightarrow\text{KNN/Clustering}
\]
\[
\text{Inner Product}\rightarrow\text{Similarity}
\]
\[
\text{Hilbert Space}\rightarrow\text{Kernel Methods}
\]

These concepts form the foundation for the linear algebra topics that follow.
\end{frame}

\section{Vectors}

\begin{frame}{Vectors}
A vector is a mathematical object having magnitude and direction.

\[
\mathbf v=
\begin{bmatrix}
v_1\\v_2\\\vdots\\v_n
\end{bmatrix}
\]

Examples:
\[
\mathbf v=\begin{bmatrix}3\\2\end{bmatrix},
\qquad
\mathbf u=\begin{bmatrix}1\\4\\2\end{bmatrix}
\]

Magnitude:
\[
\|\mathbf v\|=\sqrt{v_1^2+v_2^2+\cdots+v_n^2}
\]

In machine learning, a data point can be represented as a feature vector:
\[
\mathbf x=[x_1,x_2,\ldots,x_n]^T.
\]
\end{frame}

\begin{frame}{Coordinate System}
A coordinate system provides a way to represent points and vectors.

\textbf{2D:}
\[
P=(x,y),\qquad
\mathbf p=\begin{bmatrix}x\\y\end{bmatrix}
\]

\textbf{3D:}
\[
P=(x,y,z),\qquad
\mathbf p=
\begin{bmatrix}x\\y\\z\end{bmatrix}
\]

Standard basis vectors in $\mathbb R^3$:
\[
\mathbf i=\begin{bmatrix}1\\0\\0\end{bmatrix},
\quad
\mathbf j=\begin{bmatrix}0\\1\\0\end{bmatrix},
\quad
\mathbf k=\begin{bmatrix}0\\0\\1\end{bmatrix}
\]

Thus:
\[
\mathbf v=x\mathbf i+y\mathbf j+z\mathbf k.
\]
\end{frame}

\begin{frame}{Vector Addition}
For
\[
\mathbf u=\begin{bmatrix}u_1\\u_2\end{bmatrix},
\quad
\mathbf v=\begin{bmatrix}v_1\\v_2\end{bmatrix},
\]

\[
\mathbf u+\mathbf v=
\begin{bmatrix}
u_1+v_1\\u_2+v_2
\end{bmatrix}
\]

Example:
\[
\begin{bmatrix}2\\3\end{bmatrix}
+
\begin{bmatrix}4\\1\end{bmatrix}
=
\begin{bmatrix}6\\4\end{bmatrix}
\]

Geometrically, vector addition follows the parallelogram law.

\[
\boxed{\text{Resultant}=\text{diagonal of parallelogram}}
\]
\end{frame}

\begin{frame}{Scalar Multiplication}
For scalar $\alpha$ and vector $\mathbf v$:

\[
\alpha\mathbf v=
\begin{bmatrix}
\alpha v_1\\
\alpha v_2\\
\vdots\\
\alpha v_n
\end{bmatrix}
\]

Example:
\[
3\begin{bmatrix}2\\1\end{bmatrix}
=
\begin{bmatrix}6\\3\end{bmatrix}
\]

Magnitude:
\[
\|\alpha\mathbf v\|=|\alpha|\|\mathbf v\|
\]

\begin{itemize}
\item $\alpha>1$: vector becomes longer
\item $0<\alpha<1$: vector becomes shorter
\item $\alpha<0$: direction reverses
\item $\alpha=0$: zero vector
\end{itemize}
\end{frame}

\section{Linear Combinations and Basis}

\begin{frame}{Linear Combinations}
A linear combination of vectors $\mathbf v_1,\ldots,\mathbf v_n$ is

\[
c_1\mathbf v_1+c_2\mathbf v_2+\cdots+c_n\mathbf v_n.
\]

Example:
\[
\mathbf v_1=\begin{bmatrix}1\\0\end{bmatrix},
\quad
\mathbf v_2=\begin{bmatrix}0\\1\end{bmatrix}
\]

\[
3\mathbf v_1+2\mathbf v_2
=
\begin{bmatrix}3\\2\end{bmatrix}
\]

Therefore every vector in $\mathbb R^2$ can be expressed as a linear combination of the standard basis vectors.
\end{frame}

\begin{frame}{Span}
The span of vectors is the set of all their linear combinations.

\[
\operatorname{span}\{\mathbf v_1,\ldots,\mathbf v_n\}
=
\left\{
\sum_{i=1}^{n}c_i\mathbf v_i:c_i\in\mathbb R
\right\}
\]

For
\[
\mathbf e_1=\begin{bmatrix}1\\0\end{bmatrix},
\quad
\mathbf e_2=\begin{bmatrix}0\\1\end{bmatrix},
\]

\[
\operatorname{span}\{\mathbf e_1,\mathbf e_2\}
=\mathbb R^2.
\]

Geometrically:
\[
\text{one independent vector}\rightarrow\text{line}
\]
\[
\text{two independent vectors in }\mathbb R^2\rightarrow\text{plane}.
\]
\end{frame}

\begin{frame}{Basis Vectors}
A basis is a set of vectors that:

\begin{enumerate}
\item spans the vector space
\item is linearly independent
\end{enumerate}

Standard basis of $\mathbb R^2$:
\[
\mathbf e_1=\begin{bmatrix}1\\0\end{bmatrix},
\quad
\mathbf e_2=\begin{bmatrix}0\\1\end{bmatrix}
\]

Every vector has a unique representation:
\[
\mathbf v=x\mathbf e_1+y\mathbf e_2.
\]

For example:
\[
\begin{bmatrix}4\\3\end{bmatrix}
=
4\mathbf e_1+3\mathbf e_2.
\]
\end{frame}

\section{Matrices and Transformations}

\begin{frame}{Matrices as Linear Transformations}
A matrix can represent a linear transformation.

\[
A:\mathbb R^n\rightarrow\mathbb R^m
\]

\[
\mathbf x\mapsto A\mathbf x
\]

Example:
\[
A=
\begin{bmatrix}
2&0\\
0&3
\end{bmatrix},
\quad
\mathbf x=
\begin{bmatrix}x\\y\end{bmatrix}
\]

\[
A\mathbf x=
\begin{bmatrix}
2x\\3y
\end{bmatrix}
\]

The transformation scales the $x$-direction by $2$ and the $y$-direction by $3$.
\end{frame}

\begin{frame}{Matrix Multiplication as Composition}
Suppose
\[
\mathbf x\xrightarrow{B}B\mathbf x
\xrightarrow{A}A(B\mathbf x).
\]

Therefore:
\[
A(B\mathbf x)=(AB)\mathbf x.
\]

Hence:
\[
\boxed{AB=\text{composition of }B\text{ followed by }A}
\]

Matrix multiplication is generally not commutative:
\[
AB\neq BA.
\]

This reflects the fact that applying transformations in different orders generally gives different results.
\end{frame}

\begin{frame}{Three-Dimensional Linear Transformations}
A linear transformation in 3D maps:
\[
\mathbb R^3\rightarrow\mathbb R^3.
\]

Common transformations:

\begin{itemize}
\item Scaling
\item Rotation
\item Reflection
\item Shearing
\end{itemize}

Scaling:
\[
A=
\begin{bmatrix}
2&0&0\\
0&3&0\\
0&0&1
\end{bmatrix}
\]

\[
A
\begin{bmatrix}x\\y\\z\end{bmatrix}
=
\begin{bmatrix}2x\\3y\\z\end{bmatrix}.
\]
\end{frame}

\begin{frame}{3D Rotation}
Rotation around the $z$-axis:

\[
R_z(\theta)=
\begin{bmatrix}
\cos\theta&-\sin\theta&0\\
\sin\theta&\cos\theta&0\\
0&0&1
\end{bmatrix}
\]

The transformed vector is:

\[
\mathbf x'=R_z(\theta)\mathbf x.
\]

Applications:
\begin{itemize}
\item Computer graphics
\item Computer vision
\item Robotics
\item 3D image processing
\end{itemize}
\end{frame}

\section{Determinants and Inverses}

\begin{frame}{Determinant}
For
\[
A=
\begin{bmatrix}
a&b\\c&d
\end{bmatrix},
\]

\[
\boxed{\det(A)=ad-bc}
\]

Example:
\[
A=
\begin{bmatrix}
2&3\\1&4
\end{bmatrix}
\]

\[
\det(A)=2(4)-3(1)=5.
\]

Geometric interpretation:
\[
|\det(A)|=\text{area scaling factor}
\]

For a 3D transformation:
\[
|\det(A)|=\text{volume scaling factor}.
\]

If
\[
\det(A)=0,
\]
the transformation collapses dimensions and $A$ is singular.
\end{frame}

\begin{frame}{Inverse Matrices}
The inverse reverses a matrix transformation:

\[
A^{-1}A=AA^{-1}=I.
\]

For
\[
A=
\begin{bmatrix}a&b\\c&d\end{bmatrix},
\]

\[
A^{-1}
=
\frac{1}{ad-bc}
\begin{bmatrix}
d&-b\\-c&a
\end{bmatrix}
\]

provided:
\[
ad-bc\neq0.
\]

Therefore:
\[
\boxed{\det(A)\neq0\iff A^{-1}\text{ exists}}
\]
for a square matrix.
\end{frame}

\section{Column and Null Spaces}

\begin{frame}{Column Space}
The column space is the span of the columns of a matrix.

\[
\operatorname{Col}(A)
=
\operatorname{span}\{\text{columns of }A\}
\]

Example:
\[
A=
\begin{bmatrix}
1&2\\
2&4
\end{bmatrix}
\]

Since
\[
\begin{bmatrix}2\\4\end{bmatrix}
=
2\begin{bmatrix}1\\2\end{bmatrix},
\]

\[
\operatorname{Col}(A)
=
\operatorname{span}
\left\{
\begin{bmatrix}1\\2\end{bmatrix}
\right\}.
\]

The column space represents all possible outputs $A\mathbf x$.
\end{frame}

\begin{frame}{Null Space}
The null space contains all vectors mapped to the zero vector.

\[
\boxed{
\operatorname{Null}(A)=
\{\mathbf x:A\mathbf x=\mathbf0\}
}
\]

For
\[
A=
\begin{bmatrix}
1&2\\
2&4
\end{bmatrix},
\]

solve:
\[
x+2y=0.
\]

Thus:
\[
x=-2y
\]

and
\[
\mathbf x
=
t\begin{bmatrix}-2\\1\end{bmatrix}.
\]

Therefore:
\[
\boxed{
\operatorname{Null}(A)=
\operatorname{span}
\left\{
\begin{bmatrix}-2\\1\end{bmatrix}
\right\}
}
\]
\end{frame}

\section{Nonsquare Matrices}

\begin{frame}{Nonsquare Matrices as Transformations}
A matrix need not be square.

For example:
\[
A=
\begin{bmatrix}
1&0&2\\
0&1&3
\end{bmatrix}
\]

is a $2\times3$ matrix.

Therefore:
\[
\boxed{A:\mathbb R^3\rightarrow\mathbb R^2}
\]

and
\[
A
\begin{bmatrix}x\\y\\z\end{bmatrix}
=
\begin{bmatrix}
x+2z\\
y+3z
\end{bmatrix}.
\]

In general:
\[
\boxed{
m\times n\text{ matrix}:
\mathbb R^n\rightarrow\mathbb R^m
}
\]

This is useful for dimensionality reduction and feature transformations.
\end{frame}

\section{Dot Product and Duality}

\begin{frame}{Dot Product}
For
\[
\mathbf u=
\begin{bmatrix}u_1\\\vdots\\u_n\end{bmatrix},
\quad
\mathbf v=
\begin{bmatrix}v_1\\\vdots\\v_n\end{bmatrix},
\]

\[
\boxed{
\mathbf u\cdot\mathbf v
=
\sum_{i=1}^{n}u_iv_i
}
\]

or:
\[
\mathbf u\cdot\mathbf v=\mathbf u^T\mathbf v.
\]

Geometric form:
\[
\boxed{
\mathbf u\cdot\mathbf v
=
\|\mathbf u\|\|\mathbf v\|\cos\theta
}
\]

If:
\[
\mathbf u\cdot\mathbf v=0,
\]
the vectors are orthogonal.
\end{frame}

\begin{frame}{Dot Product and Duality}
A vector can define a linear functional.

For
\[
\mathbf w=
\begin{bmatrix}w_1\\w_2\end{bmatrix},
\]

define:
\[
f(\mathbf x)=\mathbf w^T\mathbf x.
\]

Thus:
\[
f(x,y)=w_1x+w_2y.
\]

The set of linear functionals on a vector space forms its dual space.

\[
\boxed{
\mathbf w\longrightarrow
\mathbf w^T(\cdot)
}
\]

This concept is important in linear algebra, optimization and machine learning.
\end{frame}

\section{Cross Product}

\begin{frame}{Cross Product}
For vectors in $\mathbb R^3$:

\[
\mathbf u=
\begin{bmatrix}u_1\\u_2\\u_3\end{bmatrix},
\quad
\mathbf v=
\begin{bmatrix}v_1\\v_2\\v_3\end{bmatrix},
\]

\[
\mathbf u\times\mathbf v=
\begin{bmatrix}
u_2v_3-u_3v_2\\
u_3v_1-u_1v_3\\
u_1v_2-u_2v_1
\end{bmatrix}.
\]

The result is perpendicular to both $\mathbf u$ and $\mathbf v$.

\[
\boxed{
(\mathbf u\times\mathbf v)\cdot\mathbf u=0
}
\]

\[
\boxed{
(\mathbf u\times\mathbf v)\cdot\mathbf v=0
}
\]
\end{frame}

\begin{frame}{Cross Product: Geometric Meaning}
Magnitude:

\[
\boxed{
\|\mathbf u\times\mathbf v\|
=
\|\mathbf u\|\|\mathbf v\|\sin\theta
}
\]

This equals the area of the parallelogram formed by $\mathbf u$ and $\mathbf v$.

\[
\boxed{
A_{\text{parallelogram}}
=
\|\mathbf u\times\mathbf v\|
}
\]

Triangle area:
\[
\boxed{
A_{\text{triangle}}
=
\frac12\|\mathbf u\times\mathbf v\|
}
\]

The cross product therefore provides a normal vector and measures area.
\end{frame}

\begin{frame}{Cross Product and Linear Transformations}
Let:
\[
\mathbf n=\mathbf u\times\mathbf v.
\]

Then $\mathbf n$ is normal to the plane containing $\mathbf u$ and $\mathbf v$.

A plane can be written as:
\[
\mathbf n^T(\mathbf x-\mathbf x_0)=0.
\]

Equivalently:
\[
\boxed{\mathbf n^T\mathbf x=c}
\]

For an invertible $3\times3$ matrix:
\[
\boxed{
(A\mathbf u)\times(A\mathbf v)
=
\det(A)A^{-T}(\mathbf u\times\mathbf v)
}
\]

This connects cross products with determinants, inverse transformations, area and orientation.
\end{frame}

\section{Cramer's Rule}

\begin{frame}{Cramer's Rule}
For:
\[
A\mathbf x=\mathbf b,
\]

Cramer's rule uses determinants to solve the system.

For:
\[
a_1x+b_1y=c_1,
\qquad
a_2x+b_2y=c_2,
\]

\[
D=
\begin{vmatrix}
a_1&b_1\\
a_2&b_2
\end{vmatrix}
\]

\[
x=\frac{D_x}{D},
\qquad
y=\frac{D_y}{D}.
\]

The determinant $D_x$ replaces the first column of $A$ by $\mathbf b$, while $D_y$ replaces the second column.
\end{frame}

\begin{frame}{Cramer's Rule: Example}
Solve:
\[
2x+y=5
\]
\[
x-y=1.
\]

\[
D=
\begin{vmatrix}
2&1\\1&-1
\end{vmatrix}
=-3
\]

\[
D_x=
\begin{vmatrix}
5&1\\1&-1
\end{vmatrix}
=-6
\]

\[
x=\frac{-6}{-3}=2.
\]

\[
D_y=
\begin{vmatrix}
2&5\\1&1
\end{vmatrix}
=-3
\]

\[
y=\frac{-3}{-3}=1.
\]

\[
\boxed{x=2,\qquad y=1}
\]
\end{frame}

\section{Change of Basis}

\begin{frame}{Change of Basis}
The same vector can have different coordinate representations in different bases.

Standard basis:
\[
E=
\left\{
\begin{bmatrix}1\\0\end{bmatrix},
\begin{bmatrix}0\\1\end{bmatrix}
\right\}
\]

Another basis:
\[
B=
\left\{
\begin{bmatrix}1\\1\end{bmatrix},
\begin{bmatrix}1\\-1\end{bmatrix}
\right\}.
\]

For:
\[
\mathbf v=
\begin{bmatrix}3\\1\end{bmatrix},
\]

we find:
\[
\mathbf v=
2\begin{bmatrix}1\\1\end{bmatrix}
+
1\begin{bmatrix}1\\-1\end{bmatrix}.
\]

Thus:
\[
[\mathbf v]_B=
\begin{bmatrix}2\\1\end{bmatrix}.
\]
\end{frame}

\begin{frame}{Change-of-Basis Matrix}
Put the basis vectors into the columns of $P$:

\[
P=
\begin{bmatrix}
1&1\\
1&-1
\end{bmatrix}.
\]

Then:
\[
[\mathbf v]_E=P[\mathbf v]_B.
\]

Therefore:
\[
\boxed{
[\mathbf v]_B=P^{-1}[\mathbf v]_E
}
\]

The vector remains the same geometrically; only its coordinate representation changes.
\end{frame}

\section{Eigenvalues and Eigenvectors}

\begin{frame}{Eigenvectors and Eigenvalues}
An eigenvector is a vector whose direction remains unchanged under a linear transformation.

\[
\boxed{
A\mathbf v=\lambda\mathbf v
}
\]

where:

\begin{itemize}
\item $\mathbf v$ = eigenvector
\item $\lambda$ = eigenvalue
\end{itemize}

The transformation only scales the eigenvector.

If $\lambda>1$, it expands.

If $0<\lambda<1$, it contracts.

If $\lambda<0$, its direction reverses.
\end{frame}

\begin{frame}{Finding Eigenvalues}
Starting with:
\[
A\mathbf v=\lambda\mathbf v
\]

\[
(A-\lambda I)\mathbf v=0.
\]

For a nonzero $\mathbf v$:
\[
\det(A-\lambda I)=0.
\]

This is the characteristic equation.

Example:
\[
A=
\begin{bmatrix}
2&0\\0&3
\end{bmatrix}.
\]

\[
\det(A-\lambda I)
=
(2-\lambda)(3-\lambda)=0.
\]

Hence:
\[
\boxed{\lambda_1=2,\qquad\lambda_2=3}
\]
\end{frame}

\begin{frame}{Eigenvector Example}
For:
\[
A=
\begin{bmatrix}
2&0\\0&3
\end{bmatrix},
\]

for $\lambda=2$:

\[
A
\begin{bmatrix}1\\0\end{bmatrix}
=
2
\begin{bmatrix}1\\0\end{bmatrix}.
\]

Thus:
\[
\boxed{
\mathbf v_1=
\begin{bmatrix}1\\0\end{bmatrix}
}
\]

For $\lambda=3$:

\[
A
\begin{bmatrix}0\\1\end{bmatrix}
=
3
\begin{bmatrix}0\\1\end{bmatrix}.
\]

Thus:
\[
\boxed{
\mathbf v_2=
\begin{bmatrix}0\\1\end{bmatrix}
}
\]
\end{frame}

\begin{frame}{Applications of Eigenvalues and Eigenvectors}
Eigenvalues and eigenvectors are used in:

\begin{itemize}
\item Principal Component Analysis (PCA)
\item Image processing
\item Dimensionality reduction
\item Dynamical systems
\item Differential equations
\item Google PageRank
\item Quantum mechanics
\end{itemize}

In PCA, eigenvectors of the covariance matrix identify important directions of variation in the data.
\end{frame}

\section{Vector Spaces}

\begin{frame}{Vector Spaces}
A vector space is a set closed under vector addition and scalar multiplication and satisfying the vector-space axioms.

Examples:
\[
\mathbb R^2,\qquad
\mathbb R^3,\qquad
\mathbb R^n
\]

Polynomial vector space:
\[
P_2=
\{a+bx+cx^2:a,b,c\in\mathbb R\}.
\]

A polynomial:
\[
p(x)=a+bx+cx^2
\]

can be represented as:
\[
\begin{bmatrix}a\\b\\c\end{bmatrix}.
\]

Therefore:
\[
\boxed{P_2\cong\mathbb R^3}
\]
\end{frame}

\section{Machine Learning Connection}

\begin{frame}{Vectors and the Perceptron}
A perceptron receives an input vector:

\[
\mathbf x=
\begin{bmatrix}
x_1\\x_2\\\vdots\\x_n
\end{bmatrix}
\]

and has a weight vector:

\[
\mathbf w=
\begin{bmatrix}
w_1\\w_2\\\vdots\\w_n
\end{bmatrix}.
\]

The weighted sum is:

\[
z=\mathbf w^T\mathbf x+b.
\]

The decision boundary is:

\[
\boxed{
\mathbf w^T\mathbf x+b=0
}
\]

which is a hyperplane.
\end{frame}

\begin{frame}{Perceptron and Angle Between Vectors}
Using the dot product:

\[
\mathbf w^T\mathbf x
=
\|\mathbf w\|\|\mathbf x\|\cos\theta.
\]

Therefore:

\[
z=
\|\mathbf w\|\|\mathbf x\|\cos\theta+b.
\]

Ignoring $b$:

\[
\theta<90^\circ
\Rightarrow
\mathbf w^T\mathbf x>0
\]

\[
\theta=90^\circ
\Rightarrow
\mathbf w^T\mathbf x=0
\]

\[
\theta>90^\circ
\Rightarrow
\mathbf w^T\mathbf x<0.
\]

Thus the angle between the input vector and weight vector influences classification.
\end{frame}

\begin{frame}{Complete Conceptual Connection}
\[
\boxed{
\text{Vectors}
\rightarrow
\text{Linear Combinations}
\rightarrow
\text{Span}
\rightarrow
\text{Basis}
\rightarrow
\text{Vector Spaces}
}
\]

\[
\boxed{
\text{Matrices}
\rightarrow
\text{Linear Transformations}
\rightarrow
\text{Composition}
\rightarrow
\det(A)
\rightarrow
A^{-1}
}
\]

\[
\boxed{
\text{Dot Product}
\rightarrow
\text{Angle}
\rightarrow
\text{Orthogonality}
\rightarrow
\text{Duality}
}
\]

\[
\boxed{
\text{Eigenvectors}
\rightarrow
\text{Special Directions}
\rightarrow
\text{PCA and ML}
}
\]

\[
\boxed{
\mathbf w^T\mathbf x+b=0
\rightarrow
\text{Hyperplane}
\rightarrow
\text{Perceptron}
}
\]
\end{frame}

\begin{frame}{Summary}
\begin{itemize}
\item Vectors represent quantities with magnitude and direction.
\item Linear combinations generate spans and lead to basis vectors.
\item Matrices represent linear transformations.
\item Matrix multiplication represents composition of transformations.
\item Determinants measure area or volume scaling and determine invertibility.
\item Column space describes possible outputs; null space describes vectors mapped to zero.
\item Dot products measure alignment and provide angles and orthogonality.
\item Cross products produce normal vectors and measure area in $\mathbb R^3$.
\item Cramer's rule solves linear systems using determinants.
\item Change of basis changes coordinate representation without changing the vector itself.
\item Eigenvectors identify directions preserved by a transformation.
\item Vector spaces provide the mathematical framework for all these concepts.
\end{itemize}
\end{frame}

\begin{frame}
\centering
\Huge Thank You
\end{frame}

\end{document}
